\paragraph{临界成功相图、六倍窗半尺度缺口、局部 Cram\'er 展开与 escort 截断律}
在 bin-fold 临界资源窗的双折点极限、六倍窗零点块的半阶密度律、真实输入 \(40\)-状态核输出位势的局部 Legendre 正规形，以及 escort 温度族的两点极限已经固定之后，还可进一步抽取出下列七条彼此兼容的后果。它们分别刻画临界 bit 尺度下最优反演成功率的精确三相相图、给定目标成功率时连续预算与整数预算的反演公式、六倍窗零点块的 Lucas 半尺度律、碰撞缺口与熵缺口共享的算术半尺度缺口、典型输出密度附近的四阶 Cram\'er--Legendre 展开及其左右不对称性，以及 escort 温度族在固定双尺度 \(\{\varphi^{-1},1\}\) 上的截断极限。

\begin{theorem}[临界 bit 尺度下最优反演成功率的精确三相相图]\label{thm:xi-time-part9zk-critical-bit-success-phase-diagram}
定义 bin-fold 的最优成功率
\[
\mathrm{Succ}_m(B):=\mathrm{Succ}_m^{\mathrm{bin}}(B)
=
2^{-m}C_m^{\mathrm{bin}}(B)
\qquad (B\in \ZZ_{\ge 0}).
\]
若整数序列 \((B_m)_{m\ge 1}\) 满足
\[
s_m:=2^{B_m}\left(\frac{\varphi}{2}\right)^m\longrightarrow s\in(0,\infty),
\]
则存在极限
\[
\mathrm{Succ}_m(B_m)\longrightarrow \mathcal S_\infty(s),
\]
其中
\[
\mathcal S_\infty(s)
:=
\frac{\varphi}{\sqrt5}\min\{1,s\}
\;+\;
\frac{1}{\sqrt5}\min\{\varphi^{-1},s\}.
\]
等价地，
\[
\mathcal S_\infty(s)=
\begin{cases}
\dfrac{\varphi^2}{\sqrt5}\,s,
& 0<s\le \varphi^{-1},\\[8pt]
\dfrac{\varphi}{\sqrt5}\,s+\dfrac{1}{\varphi\sqrt5},
& \varphi^{-1}\le s\le 1,\\[10pt]
1,
& s\ge 1.
\end{cases}
\]
因此，在临界位率
\[
m\log_2\!\left(\frac{2}{\varphi}\right)
\]
的 \(O(1)\)-窗口内，最优恢复并不表现为模糊过渡，而是由
\[
s=\varphi^{-1},
\qquad
s=1
\]
两处刚性折点所控制的三相极限曲线。
\end{theorem}
\begin{proof}
由推论 \ref{cor:xi-foldbin-critical-scale-success-curve-inversion}，定义
\[
\mathcal S_m(s):=
2^{-m}\,\mathcal C_m^{\mathrm{bin,cont}}\!\left(
s\left(\frac{2}{\varphi}\right)^m
\right),
\]
则 \(\mathcal S_m\) 在每个紧集 \(K\Subset(0,\infty)\) 上一致收敛到
\[
\mathcal S_\infty(s)
=
\frac{\varphi}{\sqrt5}\min\{1,s\}
\;+\;
\frac{1}{\sqrt5}\min\{\varphi^{-1},s\}.
\]
另一方面，由定义
\[
\mathrm{Succ}_m(B_m)
=
2^{-m}C_m^{\mathrm{bin}}(B_m)
=
2^{-m}\,\mathcal C_m^{\mathrm{bin,cont}}(2^{B_m})
=
\mathcal S_m(s_m).
\]
取紧区间 \(K\Subset(0,\infty)\) 使得 \(s\in K\) 且对充分大的 \(m\) 都有 \(s_m\in K\)。于是
\[
\bigl|\mathrm{Succ}_m(B_m)-\mathcal S_\infty(s)\bigr|
\le
\bigl|\mathcal S_m(s_m)-\mathcal S_\infty(s_m)\bigr|
+
\bigl|\mathcal S_\infty(s_m)-\mathcal S_\infty(s)\bigr|.
\]
第一项由 \(\mathcal S_m\to \mathcal S_\infty\) 在 \(K\) 上的一致收敛趋于 \(0\)，第二项由 \(\mathcal S_\infty\) 的连续性和 \(s_m\to s\) 也趋于 \(0\)。故
\[
\mathrm{Succ}_m(B_m)\to \mathcal S_\infty(s).
\]
将 \(\mathcal S_\infty\) 逐段展开，即得所示分段公式。证毕。
\end{proof}

\begin{theorem}[给定目标成功率的连续临界预算反演公式]\label{thm:xi-time-part9zk-continuous-critical-budget-inversion}
对任意 \(p\in(0,1)\)，定义连续临界预算
\[
b_m^\star(p)
:=
\inf\left\{
b\in [0,\infty):\
2^{-m}\,\mathcal C_m^{\mathrm{bin,cont}}(2^b)\ge p
\right\}.
\]
则有渐近展开
\[
b_m^\star(p)
=
m\log_2\!\left(\frac{2}{\varphi}\right)
+
\log_2 s^\star(p)
+
o(1),
\]
其中反演函数 \(s^\star:(0,1)\to(0,1)\) 显式为
\[
s^\star(p)
=
\begin{cases}
\dfrac{\sqrt5}{\varphi^2}\,p,
& 0<p\le \dfrac{\varphi}{\sqrt5},\\[8pt]
\dfrac{\sqrt5\,p-\varphi^{-1}}{\varphi},
& \dfrac{\varphi}{\sqrt5}\le p<1.
\end{cases}
\]
相应的整数最小预算
\[
B_m^\star(p)
:=
\min\bigl\{B\in\ZZ_{\ge 0}:\ \mathrm{Succ}_m(B)\ge p\bigr\}
\]
满足严格等式
\[
B_m^\star(p)=\left\lceil b_m^\star(p)\right\rceil.
\]
因此，目标成功率
\[
p=\frac{\varphi}{\sqrt5}
\]
构成严格的预算转折点；其下侧与上侧分别对应两条不同斜率的线性反演律。
\end{theorem}
\begin{proof}
定义由连续临界预算所诱导的尺度变量
\[
s_m^\star(p)
:=
2^{b_m^\star(p)}\left(\frac{\varphi}{2}\right)^m.
\]
其中 \(\mathcal S_m\) 仍取自定理
\ref{thm:xi-time-part9zk-critical-bit-success-phase-diagram}
的证明。由
\[
2^{-m}\,\mathcal C_m^{\mathrm{bin,cont}}(2^b)
=
\mathcal S_m\!\left(2^b\left(\frac{\varphi}{2}\right)^m\right)
\]
可知
\[
b_m^\star(p)
=
m\log_2\!\left(\frac{2}{\varphi}\right)
+
\log_2 s_m^\star(p).
\]
又由于
\[
b\longmapsto 2^{-m}\,\mathcal C_m^{\mathrm{bin,cont}}(2^b)
\]
在 \([0,\infty)\) 上连续且单调不减，定义中的下确界必被取到，从而
\[
\mathcal S_m\bigl(s_m^\star(p)\bigr)=p.
\]

再令 \(s^\star(p)\) 为极限曲线 \(\mathcal S_\infty\) 的最小资源广义逆：
\[
s^\star(p):=
\inf\{s\ge 0:\ \mathcal S_\infty(s)\ge p\}.
\]
由定理 \ref{thm:xi-time-part9zk-critical-bit-success-phase-diagram}，\(\mathcal S_\infty\) 在 \([0,1]\) 上连续且严格递增，而 \(p\in(0,1)\) 时有 \(s^\star(p)\in(0,1)\)。

任取
\[
0<\varepsilon<\min\{s^\star(p),1-s^\star(p)\}.
\]
由严格单调性，
\[
\mathcal S_\infty(s^\star(p)-\varepsilon)<p<\mathcal S_\infty(s^\star(p)+\varepsilon).
\]
再由推论 \ref{cor:xi-foldbin-critical-scale-success-curve-inversion} 的紧集上一致收敛，可知当 \(m\) 充分大时，
\[
\mathcal S_m(s^\star(p)-\varepsilon)<p<\mathcal S_m(s^\star(p)+\varepsilon).
\]
结合 \(\mathcal S_m(s_m^\star(p))=p\) 与 \(\mathcal S_m\) 的单调性，得到
\[
s^\star(p)-\varepsilon
<
s_m^\star(p)
\le
s^\star(p)+\varepsilon.
\]
令 \(\varepsilon\downarrow 0\)，便得
\[
s_m^\star(p)\longrightarrow s^\star(p).
\]
将其代回 \(b_m^\star(p)\) 的表达式，即得
\[
b_m^\star(p)
=
m\log_2\!\left(\frac{2}{\varphi}\right)
+
\log_2 s^\star(p)
+
o(1).
\]

显式公式则由定理 \ref{thm:xi-time-part9zk-critical-bit-success-phase-diagram} 中 \(\mathcal S_\infty\) 的分段线性表达逐段求逆而得。

最后，函数
\[
b\longmapsto 2^{-m}\,\mathcal C_m^{\mathrm{bin,cont}}(2^b)
\]
在 \([0,\infty)\) 上连续且单调不减，而 \(\mathrm{Succ}_m(B)\) 正是其在整数点 \(B\in\ZZ_{\ge 0}\) 的取值。因此达到水平 \(p\) 的最小整数预算恰为连续阈值的上整化：
\[
B_m^\star(p)=\left\lceil b_m^\star(p)\right\rceil.
\]
证毕。
\end{proof}

\begin{theorem}[六倍窗零点块的精确半尺度律]\label{thm:xi-time-part9zk-sixfold-zero-block-exact-halfscale}
\leanverified{paper\_xi\_time\_part9zk\_sixfold\_zero\_block\_exact\_halfscale}
记
\[
M_m:=F_{m+2}.
\]
在六倍窗子列
\[
m+2=2d\equiv 0\pmod 6
\]
上，零点块 \(\mathcal Z_m\) 的相对规模满足精确双边界
\[
\frac{1}{L_d}
\le
\frac{|\mathcal Z_m|}{M_m}
\le
\tau_{\mathrm{div}}(m+2)\,\frac{1}{L_d},
\]
其中 \(L_d\) 为第 \(d\) 个 Lucas 数，\(\tau_{\mathrm{div}}\) 为约数函数。并且
\[
\frac{1}{L_d}
=
\varphi^{-d}+O(\varphi^{-3d})
=
\varphi^{-(m+2)/2}
+
O\!\left(\varphi^{-3(m+2)/2}\right).
\]
特别地，
\[
\frac{|\mathcal Z_m|}{M_m}
=
\exp\!\left(
-\frac{m}{2}\log\varphi+o(m)
\right),
\]
等价地，
\[
\lim_{\substack{m\to\infty\\ m+2\equiv 0\ (6)}}
\frac1m\log\frac{M_m}{|\mathcal Z_m|}
=
\frac12\log\varphi.
\]
因此，在六倍窗上，零点块并非仅仅指数稀薄，而是以由 Lucas 因子精确锁定的半尺度斜率衰减。
\end{theorem}
\begin{proof}
由推论 \ref{cor:fold-zero-half-index-multiple6}，
\[
|\mathcal Z_m|\ge F_d.
\]
而 \(M_m=F_{m+2}=F_{2d}\)，并且 Fibonacci--Lucas 恒等式给出
\[
F_{2d}=F_dL_d.
\]
因此
\[
\frac{|\mathcal Z_m|}{M_m}
\ge
\frac{F_d}{F_{2d}}
=
\frac{1}{L_d}.
\]

另一方面，由定理 \ref{thm:fold-zero-density-sparse}，
\[
|\mathcal Z_m|
\le
\tau_{\mathrm{div}}(m+2)\,F_{\lfloor (m+2)/2\rfloor}
=
\tau_{\mathrm{div}}(m+2)\,F_d.
\]
再除以 \(M_m=F_{2d}=F_dL_d\)，便得
\[
\frac{|\mathcal Z_m|}{M_m}
\le
\tau_{\mathrm{div}}(m+2)\,\frac{1}{L_d}.
\]

Lucas 数的 Binet 公式为
\[
L_d=\varphi^d+(-\varphi)^{-d},
\]
故
\[
\frac{1}{L_d}
=
\varphi^{-d}\Bigl(1+O(\varphi^{-2d})\Bigr)^{-1}
=
\varphi^{-d}+O(\varphi^{-3d}).
\]
再利用
\[
\log \tau_{\mathrm{div}}(n)=o(n),
\]
便得到
\[
\frac{|\mathcal Z_m|}{M_m}
=
\exp\!\left(-\frac{m}{2}\log\varphi+o(m)\right),
\]
以及相应的对数极限。证毕。
\end{proof}

\begin{theorem}[六倍窗上的碰撞缺口与熵缺口共享 Lucas 半尺度缺口]\label{thm:xi-time-part9zk-sixfold-notch-collision-entropy}
沿用定理 \ref{thm:xi-time-part9zk-sixfold-zero-block-exact-halfscale} 的记号，并令 \(u_m\) 为 \(\ZZ/(M_m\ZZ)\) 上的均匀分布，
\[
p_m(r):=\frac{d_m(r)}{2^m}.
\]
则在
\[
m+2=2d\equiv 0\pmod 6
\]
上，有显式下界
\[
1-\mathsf{Col}_m
\ge
\frac{1}{L_d},
\]
以及
\[
\log M_m-D_{\mathrm{KL}}(p_m\|u_m)
\ge
-\log\!\left(1-\frac{1}{L_d}\right)
\ge
\frac{1}{L_d}.
\]
特别地，两种缺口都被同一个 Lucas 半尺度
\[
\frac{1}{L_d}
=
\varphi^{-d}+O(\varphi^{-3d})
\]
从下方共同钉扎。换言之，六倍窗上的算术零点块在碰撞与熵两侧同时强制注入一个确定性的半尺度缺口。
\end{theorem}
\begin{proof}
由定理 \ref{thm:xi-time-part60ab-window6-zero-block-halfscale-information-improvement}，
\[
1-\mathsf{Col}_m
\ge
\frac{F_d}{F_{m+2}},
\qquad
\log F_{m+2}-D_{\mathrm{KL}}(p_m\|u_m)
\ge
\log\frac{F_{m+2}}{F_{m+2}-F_d}.
\]
由于 \(F_{m+2}=F_{2d}=F_dL_d\)，故
\[
\frac{F_d}{F_{m+2}}=\frac{1}{L_d},
\]
以及
\[
\log\frac{F_{m+2}}{F_{m+2}-F_d}
=
-\log\!\left(1-\frac{1}{L_d}\right).
\]
又因对 \(0\le x<1\) 恒有
\[
-\log(1-x)\ge x,
\]
遂得
\[
\log M_m-D_{\mathrm{KL}}(p_m\|u_m)
\ge
\frac{1}{L_d}.
\]
最后代入定理 \ref{thm:xi-time-part9zk-sixfold-zero-block-exact-halfscale} 中
\[
\frac{1}{L_d}=\varphi^{-d}+O(\varphi^{-3d}),
\]
即得结论。证毕。
\end{proof}

\begin{theorem}[典型输出密度附近的四阶 Cram\'er--Legendre 展开]\label{thm:xi-time-part9zk-output-cramer-legendre-fourth-order}
沿用定理 \ref{thm:xi-real-input-cramer-germ-quadratic-field-rigidity} 的记号。记
\[
P(\theta)=\log\lambda(e^\theta),
\qquad
\alpha_0:=P'(0),
\qquad
\kappa_j:=P^{(j)}(0)\ \ (j\ge 2),
\]
并令 \(I(\alpha)\) 为其 Legendre 对偶速率函数。则当 \(\delta\to 0\) 时，
\[
I(\alpha_0+\delta)
=
\frac{\delta^2}{2\kappa_2}
-\frac{\kappa_3}{6\kappa_2^3}\delta^3
+
\frac{3\kappa_3^2-\kappa_2\kappa_4}{24\kappa_2^5}\delta^4
+
O(\delta^5).
\]
将定理 \ref{thm:xi-real-input-cramer-germ-quadratic-field-rigidity} 中的闭式常数
\[
\kappa_2=\frac{1791}{200}-\frac{3983\sqrt5}{1000},
\qquad
\kappa_3=\frac{1416369}{5000}-\frac{126691\sqrt5}{1000},
\]
\[
\kappa_4=\frac{1354428303}{100000}-\frac{605718497\sqrt5}{100000}
\]
代入，可得数值展开
\[
I(\alpha_0+\delta)
=
10.2582524032149\,\delta^2
+
22.8681939097743\,\delta^3
+
72.7346495328052\,\delta^4
+
O(\delta^5).
\]
因此，Gaussian 二次主项之后的首个非平凡修正已经在三阶出现，并且四阶系数仍保持显著正值。
\end{theorem}
\begin{proof}
第一式正是定理 \ref{thm:xi-real-input-cramer-germ-quadratic-field-rigidity} 的局部展开。将其中给出的 \(\kappa_2,\kappa_3,\kappa_4\) 闭式代入
\[
\frac{1}{2\kappa_2},
\qquad
-\frac{\kappa_3}{6\kappa_2^3},
\qquad
\frac{3\kappa_3^2-\kappa_2\kappa_4}{24\kappa_2^5},
\]
直接化简即可得到所示数值系数。证毕。
\end{proof}

\begin{corollary}[典型点附近的大偏差左右不对称性]\label{cor:xi-time-part9zk-local-rate-left-right-asymmetry}
沿用定理 \ref{thm:xi-time-part9zk-output-cramer-legendre-fourth-order} 的记号。由于
\[
\kappa_3<0,
\]
故当 \(\delta\downarrow 0\) 时有
\[
I(\alpha_0+\delta)-I(\alpha_0-\delta)
=
-\frac{\kappa_3}{3\kappa_2^3}\,\delta^3
+
O(\delta^5).
\]
特别地，对充分小的 \(\delta>0\)，成立严格不等式
\[
I(\alpha_0+\delta)>I(\alpha_0-\delta).
\]
因此，在典型输出密度附近，向上偏离比同量级的向下偏离具有更大的指数代价；局部速率函数并不以 \(\alpha_0\) 为偶对称中心。
\end{corollary}
\begin{proof}
把定理 \ref{thm:xi-time-part9zk-output-cramer-legendre-fourth-order} 的展开分别代入 \(\delta\) 与 \(-\delta\)。偶次项相消，奇次项翻转，从而
\[
I(\alpha_0+\delta)-I(\alpha_0-\delta)
=
2\left(-\frac{\kappa_3}{6\kappa_2^3}\right)\delta^3
+
O(\delta^5)
=
-\frac{\kappa_3}{3\kappa_2^3}\,\delta^3
+
O(\delta^5).
\]
又因 \(\kappa_2>0\) 且 \(\kappa_3<0\)，故主导系数为正，结论成立。证毕。
\end{proof}

\begin{theorem}[escort 温度族的两尺度截断律]\label{thm:xi-time-part9zk-escort-two-scale-clip-law}
对任意 \(q\ge 0\)，令
\[
W_m^{(q)}\sim \pi_{m,q},
\]
并定义尺度化截断泛函
\[
\mathcal C_{m,q}^{\mathrm{clip}}(s)
:=
\mathbf E_{\pi_{m,q}}\!\left[
\min\!\left\{
\left(\frac{\varphi}{2}\right)^m d_m^{\mathrm{bin}}(W_m^{(q)}),\,s
\right\}
\right],
\qquad
s\ge 0.
\]
则当 \(m\to\infty\) 时，对每个固定 \(s\ge 0\) 都有
\[
\mathcal C_{m,q}^{\mathrm{clip}}(s)
\longrightarrow
\mathcal C_\infty^{(q)}(s),
\]
其中
\[
\mathcal C_\infty^{(q)}(s)
=
\bigl(1-\theta(q)\bigr)\min\{1,s\}
+
\theta(q)\min\{\varphi^{-1},s\},
\qquad
\theta(q)=\frac{1}{1+\varphi^{q+1}}.
\]
特别地，对所有温度 \(q\) 而言，折点始终固定在
\[
s=\varphi^{-1},
\qquad
s=1,
\]
而 \(q\) 仅通过 \(\theta(q)\) 改变这两条刚性尺度支之间的权重分配，而不改变其位置。
\end{theorem}
\begin{proof}
由定理 \ref{thm:xi-time-part9ob-escort-quantile-phase-diagram} 的证明，缩放随机变量
\[
Z_m^{(q)}
:=
\left(\frac{\varphi}{2}\right)^m d_m^{\mathrm{bin}}(W_m^{(q)})
\]
满足弱收敛
\[
Z_m^{(q)}\Longrightarrow Z^{(q)},
\]
其中
\[
\mathbf P\bigl(Z^{(q)}=1\bigr)=1-\theta(q),
\qquad
\mathbf P\bigl(Z^{(q)}=\varphi^{-1}\bigr)=\theta(q).
\]
对固定 \(s\ge 0\)，函数
\[
x\longmapsto \min\{x,s\}
\]
在 \([0,\infty)\) 上有界且连续。因此由弱收敛与连续映射定理，
\[
\mathcal C_{m,q}^{\mathrm{clip}}(s)
=
\mathbf E\!\left[\min\{Z_m^{(q)},s\}\right]
\longrightarrow
\mathbf E\!\left[\min\{Z^{(q)},s\}\right].
\]
而右端对该两点分布直接计算即得
\[
\mathbf E\!\left[\min\{Z^{(q)},s\}\right]
=
\bigl(1-\theta(q)\bigr)\min\{1,s\}
+
\theta(q)\min\{\varphi^{-1},s\}.
\]
证毕。
\end{proof}

\noindent
由此，critical-scale 预言机恢复、六倍窗零点块、典型密度附近的局部大偏差以及 escort 温度重加权之间又出现了一条新的共同骨架：在资源侧，成功率相图与预算反演都由 \(\{\varphi^{-1},1\}\) 两个折点完全决定；在频域侧，六倍窗零点块把碰撞缺口与熵缺口共同钉扎在 Lucas 半尺度 \(L_d^{-1}\) 上；在输出位势侧，局部速率函数的三阶项已经强制右尾代价严格大于左尾；而在 escort 侧，任意温度 \(q\) 只会重新分配两条固定尺度支的权重，而不会移动其位置。于是，阈值反演、暗带稀疏、非高斯偏斜与温度重加权四条接口便在同一组黄金双尺度与半尺度缺口律之下闭合。

\endinput
